\section*{अध्याय \ref{ch:g12:catalysis} --- उत्प्रेरण}

\begin{solution}{exo:g12:catalysis:1}
\omterm{def:g7:chemical-reactions:reactant}{अभिकारक}: \ce{H2O2}। उत्पाद: \ce{H2O} और \ce{O2}। \omterm{def:g12:catalysis:catalyst}{उत्प्रेरक}: \ce{I-}
(चरण 1 में खपता है, चरण 2 में वापस मिलता है)। मध्यवर्ती: \ce{IO-} (चरण 1 में
बनता है, चरण 2 में खपता है)।
\end{solution}

\begin{solution}{exo:g12:catalysis:2}
समांगी; विषमांगी; विषमांगी; समांगी।
\end{solution}

\begin{solution}{exo:g12:catalysis:3}
अनुत्प्रेरित रास्ता। आरंभिक और अंतिम स्तर एक ही हैं: वही
उत्पाद बनते हैं, उतनी ही ऊर्जा मुक्त होती है; केवल रास्ता भिन्न है।
\end{solution}

\begin{solution}{exo:g12:catalysis:4}
वह एक चरण में खपता है और दूसरे में वापस मिलता है: वह चरणों के योग के दोनों
ओर आता है, और कट जाता है।
\end{solution}

\begin{solution}{exo:g12:catalysis:5}
अभिक्रिया सतह पर होती है: सतह जितनी बड़ी, उतने ही अधिक
\omterm{def:g7:atoms-and-molecules:molecule}{अणु} एक साथ अभिक्रिया कर सकते हैं, (अक्सर महँगे) \omterm{def:g12:catalysis:catalyst}{उत्प्रेरक} के
उसी द्रव्यमान के लिए।
\end{solution}

\begin{solution}{exo:g12:catalysis:6}
योग: \ce{2Fe^{3+} + 2H2O2 + 2Fe^{2+} + 2H+ -> 2Fe^{2+} + O2 + 2H+ +
2Fe^{3+} + 2H2O}। आयरन \omterm{def:g9:ions:ion}{आयन} और हाइड्रोजन \omterm{def:g9:ions:ion}{आयन} दोनों
ओर आते हैं और कट जाते हैं: \ce{2H2O2 -> 2H2O + O2}।
\end{solution}

\begin{solution}{exo:g12:catalysis:7}
\qty{72}{mL} का आधा \qty{36}{mL} है: \omterm{def:g12:catalysis:catalyst}{उत्प्रेरक} के बिना लगभग \qty{35}{min} बाद
और उसके साथ \qty{3.5}{min} बाद पहुँचता है। \omterm{def:g12:catalysis:catalyst}{उत्प्रेरक} इसे दस
गुना छोटा कर देता है।
\end{solution}

\begin{solution}{exo:g12:catalysis:8}
अभिक्रिया प्लैटिनम की सतह पर होती है; सीसा उसे ढक लेता है और
वहीं रहता है, इसलिए गैसें \omterm{def:g9:periodic-table-first-look:metal}{धातु} तक नहीं पहुँच पातीं। \omterm{def:g12:catalysis:catalyst}{उत्प्रेरक}
``विषाक्त'' हो जाता है: मौजूद पर बेकार।
\end{solution}

\begin{solution}{exo:g12:catalysis:9}
लगभग \qty{37}{\celsius} तक अभिक्रिया ताप के साथ तेज़ होती है,
किसी भी अभिक्रिया की तरह। उससे ऊपर \omterm{def:g12:catalysis:enzyme}{एंज़ाइम}, जो एक प्रोटीन है, अपनी
आकृति और उसके साथ अपनी सक्रियता खो देता है; \qty{80}{\celsius} पर वह नष्ट हो जाता है।
\end{solution}

\begin{solution}{exo:g12:catalysis:10}
\ce{2CO + 2NO -> 2CO2 + N2}: C 2 = 2, O 4 = 4, N 2 = 2। कार्बन मोनोऑक्साइड
नाइट्रोजन मोनोऑक्साइड द्वारा ऑक्सीकृत होता है, जो अपचयित होता है: हर प्रदूषक
दूसरे को हटा देता है।
\end{solution}

\begin{solution}{exo:g12:catalysis:11}
$n(\ce{O2}) = 0.072 / 24.1 = \qty{2.99e-3}{mol}$;
\qty{10.0}{mL} में $n(\ce{H2O2}) = 2 \times \num{2.99e-3} = \qty{5.98e-3}{mol}$:
$c = \qty{0.598}{mol/L}$।
\end{solution}

\begin{solution}{exo:g12:catalysis:12}
\ce{C2H5OH -> C2H4 + H2O}, एक \omterm{def:g12:curly-arrows:categories}{विलोपन}; \ce{C2H5OH -> CH3CHO + H2}।
\omterm{def:g12:catalysis:selective}{वरणात्मक उत्प्रेरक} कई संभव अभिक्रियाओं में से एक को दूसरों की तुलना में कहीं
अधिक तेज़ करता है, और इस तरह उत्पाद तय करता है।
\end{solution}

\begin{solution}{exo:g12:catalysis:13}
\omterm{def:g12:catalysis:catalyst}{उत्प्रेरक} \omterm{def:g10:reaction-progress-table:system}{अंतिम अवस्था} नहीं बदलता: \omterm{def:g12:catalysis:catalyst}{उत्प्रेरक} के साथ और उसके बिना वाले वक्र
ऑक्सीजन के एक ही आयतन पर समाप्त होते हैं। उत्पाद की उतनी ही मात्रा
मिलती है, बस जल्दी।
\end{solution}

\begin{solution}{exo:g12:catalysis:14}
$\num{6.02e23} / \num{5e6} = \qty{1.2e17}{s}$, लगभग चार अरब वर्ष।
इसे एक सेकंड में करने के लिए: \num{1.2e17} \omterm{def:g12:catalysis:enzyme}{एंज़ाइम} \omterm{def:g7:atoms-and-molecules:molecule}{अणु}, अर्थात
\qty{2e-7}{mol}।
\end{solution}

\begin{solution}{exo:g12:catalysis:15}
अमोनिया \ce{NH3}: $150 \times 17.0 / 14.0 = \qty{182}{Mt}$। नाइट्रोजन
\omterm{def:g7:atoms-and-molecules:atom}{परमाणु}: $\num{1.50e14} / 14.0 = \qty{1.07e13}{mol}$, अर्थात
\ce{N2} के \qty{5.36e12}{mol}।
\end{solution}

\begin{solution}{pb:g12:catalysis:1}
\textbf{1.} \ce{2H2O2 -> 2H2O + O2}।

\textbf{2.} हाँ: हाइड्रोजन परॉक्साइड युग्म
\ce{H2O2}/\ce{H2O} का \omterm{def:g11:redox:oxidant}{ऑक्सीकारक} है और युग्म \ce{O2}/\ce{H2O2} का \omterm{def:g11:redox:oxidant}{अपचायक}; वह
स्वयं को ऑक्सीकृत और अपचयित करता है।

\textbf{3.} \omterm{def:g12:catalysis:catalyst}{उत्प्रेरक} के बिना अभिक्रिया अत्यंत धीमी है।

\textbf{4.} $20 / 22.4 = \qty{0.893}{mol}$।

\textbf{5.} प्रति \omterm{def:g10:the-mole:amount}{मोल} ऑक्सीजन हाइड्रोजन परॉक्साइड के दो \omterm{def:g10:the-mole:amount}{मोल}:
\qty{1.79}{mol}।

\textbf{6.} \qty{1.79}{mol/L}।

\textbf{7.} $M = 2 \times 1.0 + 2 \times 16.0 = \qty{34.0}{g/mol}$;
$C_m = 1.79 \times 34.0 = \qty{60.7}{g/L}$।

\textbf{8.} आधा: \qty{0.893}{mol/L}।

\textbf{9.} कैटालेज़: एक \omterm{def:g12:catalysis:enzyme}{एंज़ाइम}, घुला हुआ, एक जैविक \omterm{def:g12:catalysis:catalyst}{उत्प्रेरक};
आयरन(III) \omterm{def:g9:ions:ion}{आयन}: \omterm{def:g12:catalysis:homogeneous}{समांगी उत्प्रेरण}।

\textbf{10.} \ce{2Fe^{3+} + H2O2 -> 2Fe^{2+} + O2 + 2H+} और
\ce{2Fe^{2+} + H2O2 + 2H+ -> 2Fe^{3+} + 2H2O}; उनका योग
\ce{2H2O2 -> 2H2O + O2} है।

\textbf{11.} \omterm{def:g12:catalysis:catalyst}{उत्प्रेरक} के साथ आयतन कहीं अधिक तेज़ी से बढ़ता है; दोनों वक्र
एक ही अंतिम आयतन पर स्थिर हो जाते हैं।

\textbf{12.} आयरन(III) \omterm{def:g9:ions:ion}{आयन} (नारंगी) पहले चरण में आयरन(II) (हल्के
हरे) में बदलते हैं और दूसरे में वापस आयरन(III) में; अंत में
सारा आयरन फिर से आयरन(III) होता है।

\textbf{13.} उबालने से \omterm{def:g12:catalysis:enzyme}{एंज़ाइम} नष्ट हो गया है।

\textbf{14.} $0.100 \times 20 = \qty{2.0}{L}$।

\textbf{15.} $n(\ce{O2}) = 0.100 \times 0.893 = \qty{0.0893}{mol}$, अर्थात
$0.0893 \times 24.1 = \qty{2.15}{L}$।

\textbf{16.} धीमा अपघटन ऑक्सीजन छोड़ता है, जो कसकर बंद बोतल में
दाब बढ़ा देती।

\textbf{17.} प्रति लीटर ऑक्सीजन के $1.50 / 2 = \qty{0.750}{mol}$, अर्थात
$0.750 \times 22.4 = \qty{16.8}{L}$: ``16.8 आयतन''।

\textbf{18.} $(1.79 - 1.50)/1.79 \approx \qty{16}{\percent}$।

\textbf{19.} \qty{1.79}{mol/L}।
\end{solution}
